\begin{problem}{}{}{A Walk Through Combinatorics - Ch.9 Ex.3 p.225}
Let \(G\) be a simple graph on 10 vertices and 28 edges. Prove that \(G\)
contains a cycle of length 4.
\end{problem}

\begin{solution}
    It seems like we are having too many edges so that we can always form
    a cycle of length 4. Saying that there always exists a cycle of length 4 is equivalent to
    saying that there always exists two vertices A and B such that the intersection of their
    set of neighbors has at least two vertices. In other words, we want to find two vertices
    A and B such that there are at least two or more vertices that connect to both A and B.

    So why there always exists such two vertices A and B? The intuition here is that if A connects
    to a lot of neighbors and B also has a lot of neighbors then it's likely that we can find
    at least two vertices connecting to both. So let A and B be two vertices with the highest degrees.
    Since G has 28 edges, sum of all degrees of vertices = 56. It follows that \(deg(A) + deg(B) \ge 12\). 
    Even if there is an edge between A and B, there are still at least 10 edges between other 
    8 points and A or B. This implies at least two or more vertices that connect to both A and B.
\end{solution}